% ============================================================================
% Курсов проект, ОКС „Магистър“, ФКСТ, Технически университет, София
% Дисциплина: Надграждащи технологии за бази данни (семестър II)
%
% Тема: клиентска библиотека за достъп до бази данни, която реализира мрежовите
%       протоколи на PostgreSQL и MySQL пряко, без слой на диспечер на драйвери.
%
% Компилация:  latexmk -pdf Main.tex        (pdfLaTeX, НЕ XeLaTeX)
% Форматиране: preamble.tex  (нормите са на ФКСТ, стойностите не се пипат наслуки)
%
% Този файл е СКЕЛЕТЪТ: структурата и редът на разделите. Съдържанието живее
% в chapters/. Незапълнените места са маркирани с \TODO{...} и се намират с
%     grep -rn 'TODO' chapters/ Main.tex references.bib
% ============================================================================
\documentclass[12pt]{article}

% ============================================================================
% preamble.tex — форматиране по изискванията на ФКСТ.
%
% Всяка секция тук съответства на правило от REQUIREMENTS.md (Част А) и на
% реда му в таблицата на превода (Част Б). Не пипай стойностите без да
% провериш срещу REQUIREMENTS.md — те са нормативни, не са вкус.
%
% Двигател: pdfLaTeX (T2A + tempora). НЕ XeLaTeX — вж. REQUIREMENTS.md Б.0.
% ============================================================================

\usepackage{cmap}
\usepackage[T2A]{fontenc}
\usepackage[utf8]{inputenc}
\usepackage[bulgarian]{babel}

% --- Типография (А.4) ---------------------------------------------------
% tempora = Times-метрично съвместим шрифт с пълна кирилица под T2A.
\usepackage{tempora}
\usepackage[top=2.5cm, bottom=2.5cm, left=3.0cm, right=2.0cm, a4paper]{geometry}
\usepackage{setspace}
\setstretch{1.5}                     % „междуредие 1.5 lines“ — калибруем, вж. Б.2
% „без пренасяния на срички“ (А.4) — чрез ЗАБРАНИТЕЛНИ НАКАЗАНИЯ, не чрез
% \usepackage[none]{hyphenat}. Разликата е съществена: `hyphenat[none]` задава
% \hyphenchar=-1 и така изключва самия механизъм за прекъсване, от който зависи
% `breaklines` на listings — дългите редове в листингите тогава НЕ се пренасят и
% излизат извън полето (1227 препълвания дойдоха оттам). Наказанията постигат същия
% резултат за текста (TeX никога не пренася), но оставят механизма достъпен.
\usepackage[expansion=false]{microtype}
\DisableLigatures{encoding=*,family=tt*}
\hyphenpenalty=10000
\exhyphenpenalty=10000
\emergencystretch=1.5em              % цената на justified-без-преноси, вж. Б.4
\sloppy
\usepackage{indentfirst}             % отстъп и на ПЪРВИЯ абзац след заглавие, вж. Б.3
\setlength{\parindent}{1.25cm}
\setlength{\parskip}{0pt}

% --- Номерация на страниците: долу вдясно (А.4) -------------------------
\usepackage{fancyhdr}
\pagestyle{fancy}
\fancyhf{}
\renewcommand{\headrulewidth}{0pt}
\fancyfoot[R]{\thepage}

% --- Бележки под линия: 10 pt, single, justified (А.4) -------------------
\usepackage[hang,flushmargin]{footmisc}
\usepackage{ragged2e}
\renewcommand{\footnotelayout}{\fontsize{10}{12}\selectfont\setstretch{1.0}\justifying}

% --- Заглавия на раздели / параграфи / подпараграфи (А.7) ---------------
\usepackage{titlesec}
\newcommand{\sectionbreak}{\clearpage}   % „всеки раздел започва на нова страница“

\renewcommand{\thesection}{\Roman{section}}
\renewcommand{\thesubsection}{\arabic{section}.\arabic{subsection}}
\renewcommand{\thesubsubsection}{\arabic{section}.\arabic{subsection}.\arabic{subsubsection}}

% Раздел: римски, ALL CAPS, центриран, bold 14 pt
\titleformat{\section}[block]
  {\fontsize{14}{16.8}\bfseries\centering}{\thesection. }{0pt}{\MakeUppercase}

% Параграф: арабски, bold 14 pt, отстъп 1.25 cm
\titleformat{\subsection}[block]
  {\fontsize{14}{16.8}\bfseries}{\thesubsection. }{0pt}{}
\titlespacing*{\subsection}{1.25cm}{*3}{*2}

% Подпараграф: bold 13 pt, отстъп 2.0 cm
\titleformat{\subsubsection}[block]
  {\fontsize{13}{15.6}\bfseries}{\thesubsubsection. }{0pt}{}
\titlespacing*{\subsubsection}{2.0cm}{*2}{*1}

% --- Съдържание: разделите с ГЛАВНИ букви, 12 pt, bold (А.6) -------------
\usepackage{etoolbox}
\usepackage[titles]{tocloft}
% NB: кръпката ТРЯБВА да е СЛЕД tocloft — пакетът предефинира \l@section и
% иначе я изтрива (тихо: няма грешка, просто разделите излизат с малки букви).
\makeatletter
\patchcmd{\l@section}{#1}{\MakeUppercase{#1}}{}{}
\makeatother
\renewcommand{\cftsecleader}{\cftdotfill{\cftdotsep}}
\renewcommand{\cftsubsecleader}{\cftdotfill{\cftdotsep}}
\renewcommand{\cftsubsubsecleader}{\cftdotfill{\cftdotsep}}
\renewcommand{\cftsecfont}{\fontsize{12}{14}\bfseries}
\renewcommand{\cftsecpagefont}{\fontsize{12}{14}\bfseries}
\renewcommand{\cftsecaftersnum}{.\enspace}
\settowidth{\cftsecnumwidth}{VIII.\enspace}
\renewcommand{\cftsubsecaftersnum}{.}
\renewcommand{\cftsubsubsecaftersnum}{.}

% --- Надписи: таблици ОТГОРЕ, фигури ОТДОЛУ (А.8, А.9) ------------------
\usepackage{caption}
\usepackage{subcaption}
\DeclareCaptionLabelFormat{tablelabel}{Таблица #2}
\captionsetup[table]{labelformat=tablelabel, labelsep=period, position=above,
                     font={bf,small}, justification=centering, singlelinecheck=true}
\DeclareCaptionLabelFormat{figurelabel}{Фигура #2}
\captionsetup[figure]{labelformat=figurelabel, labelsep=period, position=below,
                      font={bf,small}, justification=centering, singlelinecheck=true}

% Препратки в текста: „(Табл. 3)“ / „(Фиг. 7)“ — задължителни по А.8/А.9.
\newcommand{\tabref}[1]{(Табл.~\ref{#1})}
\newcommand{\figref}[1]{(Фиг.~\ref{#1})}

% --- Математика ---------------------------------------------------------
\usepackage{amsmath,amssymb,amsfonts}
\usepackage{amsthm}

% --- Графики ------------------------------------------------------------
\usepackage{graphicx}
\usepackage{float}
\usepackage{tikz}
\usetikzlibrary{decorations.pathreplacing,arrows.meta,positioning,shapes.geometric,fit,calc}
\usepackage{adjustbox}
\usepackage{pdfpages}     % за подвързаните задание и декларация, вж. Б.5

% --- Таблици ------------------------------------------------------------
\usepackage{array}
% Ляво подравнена колона с фиксирана ширина. Тесните `p{}` колони с ДВУСТРАННО
% подравняване (наследено от основния текст) разкъсват редовете с огромни
% междудумни интервали — правилото за justified от А.4 се отнася за ОСНОВНИЯ текст,
% не за клетки от 3 cm. Всички таблици с текстови клетки ползват L.
\newcolumntype{L}[1]{>{\RaggedRight\arraybackslash}p{#1}}
\usepackage{booktabs}
\usepackage[table]{xcolor}% \cellcolor за топлинната карта
\usepackage{longtable}    % „продължение“ при пренасяне на таблица (А.8)
\usepackage{multirow}
\usepackage{tabularx}

% --- Помощни ------------------------------------------------------------
\usepackage{hyperref}
\usepackage{xurl}
\usepackage{enumitem}
\usepackage{textcomp}
\usepackage{listings}
\urlstyle{same}
% КРИТИЧНО: `breaklines` на listings вмъква ДИСКРЕЦИОННИ прекъсвания, а глобалните
% \hyphenpenalty/\exhyphenpenalty=10000 (правилото „без пренасяния“, А.4) ги правят
% безкрайно скъпи — тогава дългите редове в листингите НЕ се пренасят и излизат извън
% полето. Проверено изолирано: с наказанията 1 препълване, без тях 0. Наказанията се
% нулират САМО вътре в листинг, така че текстът си остава без пренасяния.
\makeatletter
\lst@AddToHook{PreInit}{\hyphenpenalty=50\exhyphenpenalty=50}
\makeatother
\renewcommand{\lstlistingname}{Листинг}
\renewcommand{\lstlistlistingname}{Списък на листингите}

% Italic САМО за ключови понятия (А.4) — семантичен макрос, за да е ревизуемо.
\newcommand{\term}[1]{\emph{#1}}
% Инлайн код/идентификатори.
\DeclareRobustCommand{\code}[1]{\texttt{#1}}
% Глифи за таблиците с присъди (генерирани от tools/gen_tex.py).
\newcommand{\cmark}{\ensuremath{\checkmark}}
\newcommand{\xmark}{\ensuremath{\times}}

% Числата от борда — генерирани, за да не може изречение да противоречи на таблица.
% Място за генерираните макроси с измерените числа. Файлът се изисква от
% preamble.tex, който е взет непроменен от шаблона на факултета и го включва
% безусловно. Празен е, защото измерване още няма: всяко число в отчета се
% появява едва след като измервателната програма го произведе.
\endinput

% Второ място за генерирани макроси, изисквано от непроменения preamble.tex.
% Празно по същата причина като board_macros.tex.
\endinput


% Заснемане с тънка черна рамка. Рамката НЕ е украса: интерфейсите със светла тема са
% почти бели и се сливат с бялата страница, така че без нея не се вижда къде свършва
% изображението и започва листът — а размерът на изображението е част от твърдението.
% \fboxsep=0pt прилепва рамката към пиксела; 0.4 pt е най-тънката, която се възпроизвежда
% надеждно при печат.
\newcommand{\capture}[2][0.44\textwidth]{%
  {\setlength{\fboxsep}{0pt}\setlength{\fboxrule}{0.4pt}%
   \fbox{\includegraphics[width=#1]{#2}}}}

% Съпоставка „еталон | система“ за един интерфейс.
%   #1 лента  #2 ключ на интерфейса  #3 тема  #4 описание
\newcommand{\evidencepair}[4]{%
  \begin{figure}[H]\centering
    \begin{tabular}{@{}cc@{}}
      \small\textbf{еталон} & \small\textbf{система} \\
      \capture{evidence/#1/maui/#2_#3.png} &
      \capture{evidence/#1/cpp/#2_#3.png} \\
    \end{tabular}
    \caption{#4}\label{fig:ev-#1-#2-#3}
  \end{figure}}

% Видим маркер за незапълнено място. Търси се с: grep -rn 'TODO' chapters/
\newcommand{\TODO}[1]{\textcolor{red}{\textbf{[TODO: #1]}}}

\tikzset{
    block/.style={draw, rounded corners, fill=blue!5, align=center, font=\small,
        minimum width=3.2cm, minimum height=1.05cm, inner sep=6pt},
    accent/.style={draw, rounded corners, fill=green!8, align=center, font=\small,
        minimum width=3.0cm, minimum height=0.95cm, inner sep=6pt},
    danger/.style={draw, rounded corners, fill=red!6, align=center, font=\small,
        minimum width=3.0cm, minimum height=0.95cm, inner sep=6pt},
    note/.style={draw, dashed, rounded corners, fill=gray!8, align=left, font=\small,
        inner sep=6pt},
    line/.style={-Latex, thick}
}

% --- Единично разстояние в нетекстови среди -----------------------------
\AtBeginEnvironment{tikzpicture}{\setstretch{1.1}}
\AtBeginEnvironment{lstlisting}{\setstretch{1.1}}
\AtBeginEnvironment{verbatim}{\setstretch{1.1}}
\AtBeginEnvironment{tabular}{\setstretch{1.1}}
\AtBeginEnvironment{longtable}{\setstretch{1.1}}

\BeforeBeginEnvironment{tikzpicture}{\begin{adjustbox}{max width=\textwidth}}
\AfterEndEnvironment{tikzpicture}{\end{adjustbox}}

% --- Оцветяване на C++ код ----------------------------------------------
\definecolor{cppblue}{RGB}{0,0,180}
\definecolor{cppgreen}{RGB}{0,128,0}
\definecolor{cppred}{RGB}{180,0,0}
\definecolor{cppgray}{RGB}{128,128,128}

\lstset{
    language=C++,
    basicstyle=\footnotesize\ttfamily,
    % Пренасяне на дълги редове. `breakatwhitespace=false` е СЪЩЕСТВЕНО: коментарите в
    % изходния текст стигат 110 знака без удобно място за прекъсване, и при `true` редът
    % изобщо не се пренася, а излиза извън полето. `linewidth` се задава изрично, защото
    % рамката и левият отстъп иначе не се приспадат от наличната ширина.
    breaklines=true, breakatwhitespace=false, breakautoindent=false, breakindent=0pt,
    linewidth=\dimexpr\linewidth-2.5em\relax,
    columns=flexible, keepspaces=true,
    postbreak=\mbox{\textcolor{cppgray}{\ensuremath{\hookrightarrow}}\space},
    xleftmargin=1.5em, xrightmargin=0.5em,
    keywordstyle=\color{cppblue}\bfseries,
    stringstyle=\color{cppred},
    commentstyle=\color{cppgreen},
    identifierstyle=\color{black},
    showstringspaces=false,
    numbers=left, numberstyle=\tiny\color{cppgray}, numbersep=8pt,
    aboveskip=1.5em, belowskip=1.5em,
    frame=single, framerule=0.5pt, rulecolor=\color{cppgray},
    tabsize=4, captionpos=b,
    morekeywords={concept, requires, constexpr, consteval, constinit, noexcept,
                  nullptr, override, final, co_await, co_return, co_yield,
                  expected, unexpected, span, string_view, unique_ptr, shared_ptr},
    emph={maui, view, element, handler, layout, page, application, bindable_object,
          view_ref, weak_ref, observable, grid, stack_layout, content_page},
    emphstyle=\color{cppblue},
    inputencoding=utf8,
    % Не-ASCII знаци, срещащи се в коментарите на изходния текст. Без тези редове
    % \lstinputlisting отказва с „Invalid UTF-8 byte sequence“ — listings не разбира
    % многобайтов вход, който не е в literate таблицата.
    literate={—}{{---}}1 {→}{{$\rightarrow$}}1 {×}{{$\times$}}1 {…}{{...}}1 {§}{{\S}}1
{а}{{\cyra}}1 {б}{{\cyrb}}1 {в}{{\cyrv}}1 {г}{{\cyrg}}1 {д}{{\cyrd}}1
             {е}{{\cyre}}1 {ж}{{\cyrzh}}1 {з}{{\cyrz}}1 {и}{{\cyri}}1 {й}{{\cyrishrt}}1
             {к}{{\cyrk}}1 {л}{{\cyrl}}1 {м}{{\cyrm}}1 {н}{{\cyrn}}1 {о}{{\cyro}}1
             {п}{{\cyrp}}1 {р}{{\cyrr}}1 {с}{{\cyrs}}1 {т}{{\cyrt}}1 {у}{{\cyru}}1
             {ф}{{\cyrf}}1 {х}{{\cyrh}}1 {ц}{{\cyrc}}1 {ч}{{\cyrch}}1 {ш}{{\cyrsh}}1
             {щ}{{\cyrshch}}1 {ъ}{{\cyrhrdsn}}1 {ь}{{\cyrsftsn}}1 {ю}{{\cyryu}}1 {я}{{\cyrya}}1
             {А}{{\CYRA}}1 {Б}{{\CYRB}}1 {В}{{\CYRV}}1 {Г}{{\CYRG}}1 {Д}{{\CYRD}}1
             {Е}{{\CYRE}}1 {Ж}{{\CYRZH}}1 {З}{{\CYRZ}}1 {И}{{\CYRI}}1 {Й}{{\CYRISHRT}}1
             {К}{{\CYRK}}1 {Л}{{\CYRL}}1 {М}{{\CYRM}}1 {Н}{{\CYRN}}1 {О}{{\CYRO}}1
             {П}{{\CYRP}}1 {Р}{{\CYRR}}1 {С}{{\CYRS}}1 {Т}{{\CYRT}}1 {У}{{\CYRU}}1
             {Ф}{{\CYRF}}1 {Х}{{\CYRH}}1 {Ц}{{\CYRC}}1 {Ч}{{\CYRCH}}1 {Ш}{{\CYRSH}}1
             {Щ}{{\CYRSHCH}}1 {Ъ}{{\CYRHRDSN}}1 {Ь}{{\CYRSFTSN}}1 {Ю}{{\CYRYU}}1 {Я}{{\CYRYA}}1,
}


\begin{document}

% ============================================================================
% Заглавна страница (А.5)
% ============================================================================
\begin{titlepage}
\thispagestyle{empty}
\begin{center}
    {\fontsize{18}{21.6}\selectfont\bfseries ТЕХНИЧЕСКИ УНИВЕРСИТЕТ, СОФИЯ\par}
    \vspace{0.5cm}
    {\fontsize{16}{19.2}\selectfont\bfseries ФАКУЛТЕТ „КОМПЮТЪРНИ СИСТЕМИ И ТЕХНОЛОГИИ“\par}

    \vspace{0.3cm}
    \begin{figure}[H]
        \centering
        \includegraphics[width=0.30\linewidth]{images/tu-logo.png}
    \end{figure}
    \vspace{0.2cm}

    {\fontsize{20}{23}\selectfont\bfseries КЛИЕНТСКА БИБЛИОТЕКА ЗА ДОСТЪП ДО БАЗИ ДАННИ ЧРЕЗ ПРЯКА РЕАЛИЗАЦИЯ НА МРЕЖОВИТЕ ПРОТОКОЛИ НА POSTGRESQL И MYSQL\par}

    \vspace{0.6cm}

    % NB: формулировката зависи от това какво възлага дисциплината. Учебният план
    % изброява и двете форми, „курсов проект“ и „курсова работа“. Ако заданието
    % казва второто, редът по-долу става „КУРСОВА РАБОТА“. Да се потвърди с
    % преподавателя преди печат.
    {\fontsize{14}{16.8}\selectfont\bfseries КУРСОВ ПРОЕКТ\par}
    \vspace{0.3cm}
    {\fontsize{14}{16.8}\selectfont\bfseries по дисциплина „Надграждащи технологии за бази данни“\par}

    \vfill

    \begin{flushleft}
    {\fontsize{14}{16.8}\selectfont\bfseries Разработил:\par}
    Алекс Цветанов, фак. № \TODO{фак. номер}\par
    Специалност: Компютърно и софтуерно инженерство, ОКС „Магистър“\par

    \vspace{0.7cm}

    {\fontsize{14}{16.8}\selectfont\bfseries Преподавател:\par}
    Доц. д-р инж. Надежда Кучмова\par
    \end{flushleft}

    \vfill

    % Годината на предаване. Да се коригира, ако проектът се предаде в друга сесия.
    {\fontsize{12}{14.4}\selectfont София, 2027\par}
\end{center}
\end{titlepage}

% ============================================================================
% Съдържание (А.3, А.6)
% ============================================================================
\newpage
\pagestyle{empty}
\setcounter{tocdepth}{2}
\tableofcontents
\newpage

\pagestyle{fancy}
\pagenumbering{arabic}

% ============================================================================
% УВОД, без номер
% ============================================================================
% УВОД, без номер. Ориентировъчен обем: 1 до 2 стр.
% Резюме на дейностите по проекта: защо е поставена задачата, каква е целта,
% какви са задачите, какъв е обхватът и как е устроено изложението.
\section*{УВОД}
\label{sec:introduction}
\addcontentsline{toc}{section}{УВОД}

Достъпът до релационна база данни от програма на C или C++ почти винаги минава през
чужд слой. Приложението извиква клиентска библиотека на доставчика, тя от своя страна
говори с процеса на сървъра по мрежов протокол, а между двете нерядко стои и диспечер
на драйвери като ODBC. Всеки слой има своя цена: собствен модел на паметта, собствено
разбиране за типовете, собствено копие на резултатните редове и собствен, обикновено
синхронен, модел на изчакване. Приложението плаща тази цена, без да я вижда, защото
слоят е даденост.

Настоящият курсов проект разглежда алтернативата. Мрежовите протоколи на PostgreSQL и
на MySQL са публично специфицирани~\cite{pg_protocol, mysql_protocol}, което означава,
че клиентът може да бъде написан срещу самата спецификация, а не срещу чужда
библиотека. Този избор премахва междинните слоеве и прави видими решения, които иначе
са скрити: къде се копира буфер, кога се сериализира стойност, колко обръщания до
сървъра струва едно подготвено заявление. Същият избор носи и очевидни рискове, които
проектът е длъжен да отчете честно, а не да премълчи.

\textbf{Целта} на проекта е да се проектира и реализира клиентска библиотека на C++20,
която реализира пряко мрежовите протоколи на PostgreSQL и MySQL, и да се измери
поведението ѝ спрямо доставчиковите клиентски библиотеки върху едно и също натоварване.

За постигането на целта са формулирани следните \textbf{задачи}:

\begin{enumerate}[topsep=0pt,itemsep=0pt,leftmargin=*]
\item преглед на стандарта SQL и на изградените върху него технологии за достъп,
      както и на спецификациите на двата разглеждани мрежови протокола;
\item анализ на възможните решения, слоест диспечер на драйвери, доставчикова
      клиентска библиотека или пряка реализация на протокола, и обосновка на избора;
\item проектиране на слоеста архитектура с ясно отделени транспорт, протоколни
      кодеци, декодер на редове, пул от връзки и слой за типово изображение;
\item програмна реализация на архитектурата за двата протокола;
\item изграждане на методика за измерване и на възпроизводима среда за нея;
\item експериментално сравнение с \code{libpq} и \code{libmysqlclient} и анализ на
      получените резултати.
\end{enumerate}

\textbf{Обхватът} на проекта е клиентската страна на комуникацията: установяване на
връзка, удостоверяване, изпълнение на прости и на подготвени заявления, декодиране на
резултатни редове и управление на пула от връзки. Извън обхвата остават
администрирането на сървъра, репликацията, разпределените транзакции и всяка форма на
съхранение от страна на клиента.

\textbf{Структура на изложението.} Раздел~\ref{sec:literature} представя предметната
област и прегледа на литературата. Раздел~\ref{sec:alternatives} излага разгледаните
алтернативи и обосновава избора. Раздел~\ref{sec:design} описва проектирането, а
Раздел~\ref{sec:implementation} програмната реализация. Раздел~\ref{sec:method}
дефинира методиката за измерване, а Раздел~\ref{sec:results} представя резултатите и
техния анализ. Заключението обобщава постигнатото и посочва посоките за продължение.


% ============================================================================
% I. Анализ на предметната област и литературен преглед
% ============================================================================
% I. Анализ на предметната област и литературен преглед.
% Задача на раздела: да покаже къде стои проектът спрямо стандарта SQL, спрямо
% изградените върху него технологии за достъп и спрямо самите мрежови протоколи.
\section{Анализ на предметната област и литературен преглед}
\label{sec:literature}

\subsection{Стандартът SQL и мястото на клиентския интерфейс в него}

Стандартът ISO/IEC 9075 описва не само езика за заявки, но и начина, по който външна
програма стига до данните. Част 1 задава общата рамка и понятията, с които останалите
части си служат~\cite{sql_framework}. Част 3 определя \term{интерфейса на ниво
извикване} (Call-Level Interface, CLI), тоест процедурен интерфейс, чрез който
приложение на общоцелеви език изпраща заявления и получава резултати, без тези
заявления да са вградени в кода при компилация~\cite{sql_cli}. Именно тази част е
шаблонът, върху който стъпват практическите технологии за достъп.

Разграничението, което стандартът прави, е съществено за целия проект. Стандартът
нормира \emph{интерфейса}, който приложението вижда, но не нормира \emph{протокола},
по който клиентът и сървърът си говорят по мрежата. Протоколът е дело на всеки
доставчик поотделно. Следователно преносимостта, която CLI обещава, се постига винаги
чрез междинен слой, който превежда стандартните извиквания към нестандартния протокол
на конкретния сървър.

\subsection{Технологии за достъп, изградени върху стандарта}

ODBC е най-разпространената реализация на идеята за CLI. Приложението се свързва с
\term{диспечер на драйвери}, който зарежда драйвер за конкретния източник на данни и
пренасочва извикванията към него~\cite{odbc_ref}. Върху този модел исторически се
надграждат OLE DB, ADO и по-късно ADO.NET, които добавят обектен модел, работа с
несвързани множества от редове и по-богата типова система~\cite{adonet}. Общото между
всички тях е броят на границите, през които минава един резултатен ред, преди да
стигне до приложението.

Отделно от релационния път стоят спецификациите за слабо структурирани данни и за
директорийни услуги. XML остава нормативната основа за обмен на документно
структурирани данни между разнородни системи~\cite{xml10}, а LDAP описва достъпа до
йерархични директорийни хранилища с протокол, който по устройство е близък до
разглежданите тук: двоичен обмен на съобщения с ясно определени етикети и дължини,
върху постоянна връзка~\cite{rfc4511}. Сравнението с LDAP е полезно, защото показва,
че прякото говорене по протокол е обичайна практика извън релационния свят.

\subsection{Мрежовите протоколи на PostgreSQL и MySQL}

Протоколът на PostgreSQL е съобщително ориентиран: всяко съобщение носи еднобайтов
етикет, дължина и тяло, а фазите на работа са ясно отделени, стартиране и
удостоверяване, прост цикъл на заявка, разширен цикъл с
\code{Parse}/\code{Bind}/\code{Execute} за подготвени заявления~\cite{pg_protocol}.
Протоколът на MySQL е организиран около пакети с последователен номер и разчита в
по-голяма степен на позиционно кодиране и на цели числа с променлива
дължина~\cite{mysql_protocol}. Двата протокола решават една и съща задача с различни
средства, което ги прави удобна двойка за общ клиентски слой.

Референтните клиентски библиотеки на двата доставчика са добре документирани и служат
за еталон в тази работа. Документацията на \code{libpq} описва подробно както
интерфейса, така и семантиката на асинхронния режим~\cite{libpq}. По-общата литература
за системи за обработка на транзакции обяснява защо клиентският слой не е неутрален
към производителността: границите на транзакциите, изчакванията и броят обръщания по
мрежата определят наблюдаваната закъснителност повече от скоростта на самия
сървър~\cite{gray1993, kleppmann2017}.

\subsection{Изводи от прегледа}

Прегледът очертава три нива, на които може да се стъпи: стандартизиран интерфейс с
диспечер на драйвери, доставчикова клиентска библиотека и самата спецификация на
протокола. Всяко следващо ниво надолу премахва слой и абстракция, но добавя
отговорност за коректност. \TODO{да се добавят два или три съпоставими академични или
индустриални източника за измервана цена на слоевете при достъп до бази данни, ако
такива бъдат намерени; ако не бъдат, това да се заяви изрично тук}


% ============================================================================
% II. Анализ на възможните решения и обосновка на избора
% ============================================================================
% II. Анализ на възможните решения и обосновка на избора.
% Тук се сравняват трите нива на достъп и се обосновава защо проектът избира
% най-ниското. Обосновката трябва да назове и цената на избора, не само ползата.
\section{Анализ на възможните решения и обосновка на избора}
\label{sec:alternatives}

\subsection{Разглеждани алтернативи}

Разгледани са три решения на една и съща задача, а именно да се изпълни заявление от
програма на C++ и да се получат резултатните редове в типизиран вид.

\textbf{Алтернатива А, слоест диспечер на драйвери.} Приложението извиква ODBC, а
диспечерът зарежда драйвер за конкретния сървър~\cite{odbc_ref}. Предимството е една
входна точка за много източници на данни, включително такива, за които никой не би
написал клиент сам. Недостатъкът е, че всеки резултатен ред минава през договорения от
стандарта тип и през буфер, зададен от приложението предварително, а асинхронната
работа не е част от модела.

\textbf{Алтернатива Б, доставчикова клиентска библиотека.} Приложението се свързва
пряко с \code{libpq} или с \code{libmysqlclient}. Отпада един слой, а семантиката е
точно тази на сървъра. Остават обаче два въпроса: библиотеките са на C и налагат
собствено управление на паметта на границата, а асинхронните им режими се различават
помежду си достатъчно, за да не се обединяват просто~\cite{libpq}.

\textbf{Алтернатива В, пряка реализация на мрежовия протокол.} Приложението кодира и
декодира съобщенията само, срещу публикуваната спецификация~\cite{pg_protocol,
mysql_protocol}. Отпадат всички междинни слоеве, буферът за получаване се притежава от
самата библиотека и декодирането може да работи върху него без копие. Цената е, че
коректността вече е отговорност на проекта, включително частите на протокола, които се
срещат рядко и затова се откриват късно.

\subsection{Критерии за сравнение}

Сравнението се води по критерии, всеки от които е проверим, а не въпрос на вкус: брой
слоеве, през които минава един ред; брой копия на данните между сокета и типизираната
стойност; наличие на естествен асинхронен модел; възможност за кеширане на подготвени
заявления от страна на клиента; обем на кода, който проектът поема да поддържа;
чувствителност към промени във версията на сървъра.

\begin{table}[H]
\centering
\caption{Съпоставка на разгледаните алтернативи по приетите критерии}
\label{tab:alternatives}
\begin{tabular}{L{4.2cm}L{3.0cm}L{3.0cm}L{3.0cm}}
\toprule
\textbf{Критерий} & \textbf{А. ODBC} & \textbf{Б. Доставчикова} & \textbf{В. Пряк протокол} \\
\midrule
Слоеве до реда        & \TODO{?} & \TODO{?} & \TODO{?} \\
Копия на данните      & \TODO{?} & \TODO{?} & \TODO{?} \\
Асинхронен модел      & \TODO{?} & \TODO{?} & \TODO{?} \\
Кеш на заявленията    & \TODO{?} & \TODO{?} & \TODO{?} \\
Поддържан код         & \TODO{?} & \TODO{?} & \TODO{?} \\
\bottomrule
\end{tabular}
\end{table}

Съпоставката по тези критерии \tabref{tab:alternatives} се попълва след като
реализацията и измерванията приключат, за да не съдържа предположения.

\subsection{Обосновка на избора}

Проектът избира алтернатива В. Основанието не е производителност сама по себе си, а
това, че само тя позволява цената на достъпа да бъде наблюдавана. При А и Б броят
копия и обръщания е скрит зад чужд код и може да се оцени само косвено. При В всяко
копие е ред в собствения изходен текст и следователно може да бъде премахнато или
защитено с аргумент. Дисциплината изисква и сравнение между слоестия модел и прякото
говорене по протокол, което тази постановка прави естествено.

Изборът се приема заедно с цената му. Проектът поема поддръжката на две независими
кодиращи схеми, поема риска от разминаване с бъдещи версии на сървърите и се отказва
от преносимостта към източници на данни извън двата разглеждани. Тези ограничения са
записани изрично, за да бъдат отчетени в заключението, а не преоткрити там.


% ============================================================================
% III. Проектиране на системата
% ============================================================================
% III. Проектиране на системата.
% Слоеве, отговорности, граници между тях. Фигурите се добавят, когато схемите
% бъдат начертани; дотогава мястото им е маркирано с \TODO.
\section{Проектиране на системата}
\label{sec:design}

\subsection{Слоеве и отговорности}

Системата е разделена на пет слоя с еднопосочна зависимост, отдолу нагоре. Най-отдолу
стои \term{транспортът}, който притежава сокета и буферите и не знае нищо за
съдържанието на байтовете. Над него стои \term{протоколният кодек}, по един за
PostgreSQL и за MySQL, който превръща поток от байтове в съобщения и обратно. Следва
\term{декодерът на редове}, който превръща тялото на едно съобщение с данни в изгледи
към стойности, без да ги копира. Над него е \term{пулът от връзки}, който управлява
живота на връзките и кеша от подготвени заявления. Най-отгоре стои \term{слоят за
типово изображение}, който свързва стойностите от протокола с типовете на C++.

Границата между всеки два слоя е тесен договор, а не общ обект със състояние. Целта е
всеки слой да се тества самостоятелно: кодекът да се захранва със записан байтов
поток, декодерът с готово тяло на съобщение, пулът с фиктивен транспорт. Това
разделение е и условието за смисленото сравнение с доставчиковите библиотеки, защото
позволява цената да се отнесе към конкретен слой.

\subsection{Протоколни кодеци}

Кодекът на PostgreSQL следва съобщителния модел на протокола: етикет, дължина, тяло, и
двата цикъла на изпълнение, простия и разширения~\cite{pg_protocol}. Кодекът на MySQL
следва пакетния модел с номер на последователност и цели числа с променлива
дължина~\cite{mysql_protocol}. Двата кодека нямат общ базов клас. Общото между тях е
изразено като изисквания към типа, тоест като \code{concept}, което позволява
останалите слоеве да са шаблонни и да не плащат за динамично изпращане.

\TODO{фигура: диаграма на слоевете и на потока от данни между сокета и типизираната
стойност; да се начертае с TikZ и да се озаглави по правилото за фигури}

\subsection{Декодиране без копие}

Резултатният ред пристига като част от получаващия буфер, притежаван от транспорта.
Декодерът не изгражда нови обекти за стойностите, а издава изгледи от вида
\code{std::span<const std::byte>} и \code{std::string\_view} към този буфер. Копие се
прави само когато приложението изрично поиска притежание на стойността. Условието за
коректност е ясно: изгледите са валидни, докато редът е жив, и това
ограничение е част от публичния договор, а не подразбираща се подробност.

Тази схема прехвърля отговорност върху приложението и затова е рискова. Проектът я
приема, но я подсигурява с проверки при активирани диагностики и с тестове, които
целенасочено използват изглед към вече невалиден ред, за да се убедят, че диагностиката го
улавя. \TODO{да се опише точният механизъм за проверка след като бъде избран}

\subsection{Пул от връзки и кеш на подготвените заявления}

Пулът е асинхронен. Заявка за връзка, която не може да бъде удовлетворена веднага,
спира изпълнението на извикващата съпрограма, вместо да блокира нишка. Всяка връзка
носи собствен кеш на подготвените заявления, защото подготвеното заявление в двата
разглеждани протокола е свързано със сесията, а не с базата. Кешът е с ограничен
размер и с политика на изместване. \TODO{избор на политика и на подразбиращ се размер,
след като измерванията покажат чувствителността към тях}

\subsection{Слой за типово изображение}

Слоят превръща стойност от протокола в типизирана стойност на C++ и обратно, като
съответствието се задава декларативно и се проверява при компилация. Целта е
несъответствие между схемата и заявената в кода структура да бъде грешка при
компилация или ясна грешка при изпълнение, а не мълчаливо погрешно четене на байтове.
\TODO{таблица със съответствието между типовете на PostgreSQL, типовете на MySQL и
типовете на C++, включително случаите без точно съответствие}


% ============================================================================
% IV. Програмна реализация
% ============================================================================
% III. Програмна реализация.
\section{Програмна реализация}
\label{sec:implementation}

Реализацията е на C++20~\cite{iso_cpp20}: съпрограмите носят асинхронния модел, без да
раздробяват логиката на протокола между обратни извиквания; \code{concept} изразява общото
между двата кодека без виртуални функции; \code{std::span} и \code{std::string\_view} са
носителят на декодирането без копие. Сглобяването е с CMake и проектът няма нито една
задължителна външна зависимост: на чиста машина стигат компилатор за C++20 и CMake. Това има
цена, платена на три места. Тестовата
рамка е собствена, 170 реда в \code{tests/check.hpp}, свързана с CTest. Измерването е със
собствен часовник върху \code{std::chrono::steady\_clock}. Хеш функциите, от които зависят
двата обмена за удостоверяване, са в \code{src/crypto.cpp}: MD5, SHA-1 и SHA-256~\cite{fips180},
HMAC и PBKDF2, всяка проверена срещу публикуваните контролни вектори, а не срещу собствения си
изход. Единствената незадължителна зависимост е \code{libpq}, търсена с
\code{find\_package(PostgreSQL QUIET)} и използвана само от измервателната програма.

\textbf{Непълно съобщение.} Транспортът може да върне произволна част от пакет, така че
кодекът трябва да може да спре и да продължи, без да губи състояние. Решението е кодекът да
няма състояние изобщо: \code{peek\_frame} гледа буфера и връща или цяло съобщение, или празна
стойност, без да консумира нищо, а консумирането е отделна стъпка след обработката. Тестът
\code{peek\_frame\_refuses\_to\_consume\_a\_partial\_message} подава всяко възможно начало на
едно съобщение, от нула до пълната дължина, и изисква празна стойност за всяко. Четенето на
съобщение е изчакваем обект и \code{await\_ready()} вика \code{peek\_frame}, тоест в съпрограма
се влиза веднъж на четене от сокета, а не веднъж на съобщение (Раздел~\ref{sec:defects}).

\textbf{Грешка от сървъра в средата на резултатен поток.} Грешка, съобщена по правилата на
протокола, задължава клиента да изчерпи обмена до синхронизиращата точка, преди връзката да
бъде използвана отново. Реализацията записва грешката, но продължава да чете до
\code{ReadyForQuery} и чак тогава хвърля изключение. Затова връзка, върху която е изпълнено
грешно заявление, остава годна: интеграционният тест
\code{pg\_live\_server\_error\_leaves\_the\_connection\_usable} след грешка с код 42P01
изпълнява ново заявление по същата връзка.

\textbf{Заемане на връзка от пула.} Заемането се състезава с изтичане на срок. Съпрограмата
не може да изчаква едновременно свободна връзка и таймер, затова едното от двете идва като
обратно извикване, което урежда другото. Общото състояние се държи през \code{shared\_ptr}, но
извикването по срок държи \code{std::weak\_ptr}, защото не се отменя при нормално даване на
връзка и доживява до срока си (Раздел~\ref{sec:defects}).
Освобождаването предава връзката пряко на чакащия. Първата реализация я връщаше в списъка на
свободните и събуждаше чакащия, което изглежда еквивалентно. Не е: съпрограмата, която
освобождава връзката, продължава да се изпълнява преди събудения чакащ и си взема същата
връзка обратно, така че единият работник завършва целия си дял, преди другият да получи
изобщо една връзка. Поведението е записано като тест, а не като число:
\code{pool\_hands\_a\_released\_connection\_straight\_to\_the\_waiter} пуска два работника по
четири заемания при пул с размер едно и изисква ред на изпълнение \code{abababab}.

\textbf{Обезсилване на кеша.} Изместване от кеша връща изместения запис, върху който се
изпраща \code{Close}; иначе всяко изместване би оставяло по едно заявление на сървъра. Ако
сървърът отговори с код 26000, тоест не познава името на кешираното заявление, записът се
изтрива и потокът се изпълнява повторно веднъж. Това се случва след \code{DISCARD ALL} или
след смяна на сървъра зад пулър и извикващият не бива да го вижда.

\textbf{Удостоверяване.} За PostgreSQL са реализирани доверен вход, парола в чист вид, MD5 и
SCRAM-SHA-256~\cite{rfc7677}. Последният не е разширение, а необходимост: PostgreSQL 14 и
следващите го използват по подразбиране, така че без него клиентът не може да се свърже със
сървър, който не е нарочно отслабен. Обменът е проверен срещу публикувания в RFC 7677 пример:
при фиксиран клиентски случаен низ реализацията възпроизвежда буква по буква и клиентското
доказателство, и подписа на сървъра. За MySQL е реализиран \code{mysql\_native\_password},
включително отговорът на искане за смяна на приставката; \code{caching\_sha2\_password} не е
реализирана и това се съобщава с ясна грешка, а не с мълчалив отказ.

\textbf{Обработка на грешките.} Разграничението е по типове. \code{protocol\_error} означава
байтов поток извън спецификацията и е неотстранима: позицията в потока е загубена.
\code{io\_error} е отказ на транспорта. \code{server\_exception} носи полетата на грешката и е
единствената, след която връзката се използва пак. Връзка с протоколна грешка или отказ на
транспорта се отбелязва като негодна чрез флага \code{poisoned\_}, а пулът проверява
\code{is\_usable()} при връщане и при заемане; тестът
\code{pool\_discards\_a\_connection\_the\_server\_closed\_underneath\_it} убива връзката,
докато тя е заета. За очаквано отсъствие се използва \code{std::optional}, не изключение.


% ============================================================================
% V. Методика за експериментална проверка
% ============================================================================
% V. Методика за експериментална проверка.
% Разделът описва ПЛАНА на експеримента и конфигурацията на машината.
% Числата от измерването са в Раздел VI и само там.
\section{Методика за експериментална проверка}
\label{sec:method}

\subsection{Какво се проверява}

Проверката има две отделни цели, които не бива да се смесват. Първата е коректност:
реализацията трябва да чете и записва това, което спецификацията предписва, и да получава
от сървъра същите стойности, които са били записани. Втората е цена: закъснителност,
брой обръщания по мрежата и цена на декодирането. Коректността е предусловие. Измерване
върху реализация, която дава различни резултати, няма смисъл.

Коректността се проверява на две нива. Модулните тестове сравняват изхода на кодеците с
байтови последователности, съставени поле по поле от спецификацията, и с публикуваните
контролни вектори за хеш функциите. Интеграционните тестове изпълняват заявления срещу
работещи сървъри и сравняват прочетеното със записаното.

\subsection{Еталони и натоварване}

Първоначалният замисъл предвиждаше сравнение с \code{libpq}~\cite{libpq} и
\code{libmysqlclient} върху едно и също натоварване. Измервателната програма съдържа
такъв клон: при конфигуриране CMake търси \code{libpq} и, ако я намери, сглобява и
еталонно измерване със същото заявление срещу същия сървър. На машината, на която са
направени измерванията от Раздел~\ref{sec:results}, \code{libpq} не е инсталирана, което
програмата съобщава изрично и не отчита никакво сравнение. Това е записано като
липсващо, а не заобиколено с оценка.

Натоварването се състои от пет случая, всеки изпълнен от една и съща програма върху една
и съща база:

\begin{enumerate}[leftmargin=*,label=\arabic*.]
\item едно заявление за един ред през простия цикъл;
\item същото заявление през разширения цикъл, при празен кеш на подготвените заявления;
\item същото заявление през разширения цикъл, при попадение в кеша;
\item голям резултат от 5000 реда с четири колони, веднъж в текстов и веднъж в двоичен формат;
\item заемане на връзка от пул при размер 1, 2, 4 и 8 и едновременност 1, 4 и 16.
\end{enumerate}

Схемата на тестовата таблица е една и съща за двата сървъра до типовете, които всеки от
тях именува различно, и е дадена в Приложение~А.

\subsection{Среда и възпроизводимост}

Средата се вдига с Docker Compose~\cite{docker_compose}, за да са фиксирани версиите
на сървърите и началното състояние на базите. Конфигурацията, при която са направени измерванията в
Раздел~\ref{sec:results}, е дадена в Табл.~\ref{tab:machine}.

\begin{table}[H]
\centering
\caption{Конфигурация на машината и на средата при измерването}
\label{tab:machine}
\begin{tabular}{L{5.0cm}L{9.0cm}}
\toprule
\textbf{Компонент} & \textbf{Стойност} \\
\midrule
Операционна система & Windows 11 Pro N, версия 10.0.26200 \\
Процесор & 12 логически ядра \\
Компилатор & g++ 15.2.0, MinGW-w64, \code{x86\_64-posix-seh} \\
Режим на сглобяване & \code{CMAKE\_BUILD\_TYPE=Release}, генератор Ninja 1.13.2 \\
Стандарт & C++20, без разширения на компилатора \\
CMake & 4.3.2 \\
Docker & 29.4.3, Compose v5.1.3 \\
PostgreSQL & образ \code{postgres:16.4}, порт 55432, удостоверяване \code{scram-sha-256} \\
MySQL & образ \code{mysql:8.0.39}, порт 33306, приставка \code{mysql\_native\_password} \\
Мрежа & обратна връзка (\code{127.0.0.1}) към контейнерите на същата машина \\
Еталонни библиотеки & \code{libpq} и \code{libmysqlclient} не са налични на машината \\
\bottomrule
\end{tabular}
\end{table}

Обстоятелството, че сървърите работят в контейнери на същата машина, а не на отделен
хост, е съществено за тълкуването на числата и е отбелязано изрично: разсейването,
което се вижда в Раздел~\ref{sec:results}, идва в голямата си част от планировчика на
виртуализационния слой, а не от клиента.

\subsection{Брой повторения и отхвърляне}

Всеки случай от първите четири се изпълнява 520 пъти, от които първите 20 се отхвърлят
като загряващи; за големия резултат повторенията са 120 при 20 загряващи, защото едно
повторение чете 5000 реда. За заемането от пул всеки работник изпълнява 40 заемания,
след като пулът е бил предварително напълнен с едновременни заемания, така че цената на
свързването да не влезе в измерването на заемането.

Отчитат се медиана, 95-и процентил, средна стойност и извадково стандартно отклонение.
Средна стойност без разсейване не позволява да се съди дали разликата между два случая
изобщо е значима, а при разпределение с тежка дясна опашка, каквото се наблюдава тук,
медианата и средната стойност се различават съществено и двете носят различна
информация.

Суровите проби от всяко повторение се записват във файл
\code{benchmarks/results/raw\_samples.csv} при подаден флаг \code{--raw}, откъдето са
изведени всички таблици в Раздел~\ref{sec:results}. Извадка от този файл е в
Приложение~Г.

\subsection{Показатели}

Измерват се четири показателя. Закъснителност на едно логическо действие, по проценти
на разпределението. Брой протоколни съобщения за едно логическо действие. Време за
декодиране на един резултатен набор, отделено от общото време чрез отделен часовник в
самия обработчик на ред. Закъснителност на заемане на връзка от пул.

Броят съобщения не е времеви и се получава чрез броене в самия клиент, което го прави
независим от машината и затова най-удобен за сравнение между реализации. Броенето се
прави от механизма за проследяване, който е част от библиотеката и е изключен по
подразбиране.

\subsection{Условия за валидност}

Измерването се приема за валидно само ако са изпълнени следните условия: и двата случая
изпълняват едно и също заявление върху една и съща база; сървърът е един и същ процес
за всички случаи в един пуск; загряващите повторения са отхвърлени; и разсейването е
отчетено. Първите две условия се осигуряват от структурата на измервателната програма,
която прави цялата подготовка веднъж в началото на пуска и след това не пипа базата;
последните две са в кода, който събира пробите. Измерване, което не отговаря на някое
от условията, се записва като невалидно и не участва в изводите.

Едно условие от първоначалния замисъл отпада: съвпадението байт по байт с еталонна
библиотека не може да бъде проверено, защото еталонна библиотека липсва на машината.
Затова изводите в Раздел~\ref{sec:results} сравняват случаи \emph{в рамките на самата
реализация} и не твърдят нищо за отношението към \code{libpq}.


% ============================================================================
% VI. Експериментални резултати и анализ
% ============================================================================
% V. Експериментални резултати и анализ.
% Всяко число тук е измерено на машината от Табл. \ref{tab:machine}.
\section{Експериментални резултати и анализ}
\label{sec:results}

Коректността е предусловие за всичко останало и е проверена изцяло: 85 модулни случая без
сървър, 6 интеграционни срещу PostgreSQL 16.4 и 6 срещу MySQL 8.0.39, всички съвпадения, нула
разминавания. Модулните случаи включват контролните вектори на MD5, SHA-1, SHA-256,
HMAC-SHA-256, PBKDF2 и пълния обмен SCRAM-SHA-256 от RFC 7677~\cite{rfc7677}; интеграционните
проверяват, че записана стойност от тип \code{numeric}, \code{timestamp} и масив се прочита
обратно със същата стойност през двоичния формат.

\begin{table}[H]
\centering
\caption{Едно логическо действие за един ред, микросекунди, $n=500$ след 20 загряващи
повторения. Съобщенията и обръщанията са преброени от механизма за проследяване}
\label{tab:latency}
\begin{tabular}{L{3.8cm}L{1.6cm}L{1.5cm}L{1.4cm}L{1.4cm}L{1.6cm}L{1.5cm}}
\toprule
\textbf{Вид заявление} & \textbf{Медиана} & \textbf{95-и пр.} & \textbf{Средно} & \textbf{Ст. откл.} & \textbf{Съобщ.} & \textbf{Обръщ.} \\
\midrule
PostgreSQL, прост цикъл          & 612 & 700 & 625 & 39 & 5  & 1 \\
PostgreSQL, разширен, пропуск    & 636 & 733 & 651 & 45 & 12 & 1 \\
PostgreSQL, разширен, попадение  & 605 & 760 & 629 & 85 & 7  & 1 \\
MySQL, текстов протокол          & 643 & 725 & 650 & 41 & не е отчетено & не е отчетено \\
\bottomrule
\end{tabular}
\end{table}

Разсейването в този пуск е малко: стандартното отклонение е под петнадесет на сто от средната
стойност. Правилото се прилага буквално: където разликата е под едно стандартно отклонение,
извод не се прави. Първото такова място е самата Табл.~\ref{tab:latency}. Трите вида
заявление за PostgreSQL се различават с по-малко от петдесет микросекунди при стандартни
отклонения между 39 и 85, тоест измерването \emph{не} подкрепя твърдение кой от трите е
по-бърз. Причината е в средата: върху обратна връзка едно допълнително протоколно съобщение
струва почти нищо.

Десните две колони дават две наблюдения. Първо, и трите случая струват едно обръщане по
мрежата: разширеният цикъл \emph{не} плаща допълнително обръщане при първото изпълнение,
защото \code{Parse}, \code{Describe}, \code{Bind}, \code{Execute} и \code{Sync} се записват в
един буфер и се изпращат заедно. Хипотезата за две обръщания би била вярна само за клиент,
който изпраща \code{Flush} и чака отговор след \code{Parse}. Второ, кешът намалява
съобщенията от 12 на 7, тоест с 42 на сто: спестени са \code{Parse} и \code{Describe} в
изпратената посока и \code{ParseComplete}, \code{ParameterDescription} и
\code{RowDescription} в получената. Тези пет съобщения не се превръщат в измеримо време тук,
затова изводът е по-тесен от очаквания: дали спестените съобщения стават спестено време
зависи от цената на едно съобщение по мрежата, а обратната връзка е най-неблагоприятният
случай за кеша.

\begin{table}[H]
\centering
\caption{Резултат от 5000 реда и четири колони, микросекунди, $n=100$. Последната колона е
времето само в декодиращите функции, измерено с отделен часовник}
\label{tab:formats}
\begin{tabular}{L{4.2cm}L{1.9cm}L{1.7cm}L{1.7cm}L{1.7cm}L{2.4cm}}
\toprule
\textbf{Случай} & \textbf{Медиана} & \textbf{95-и пр.} & \textbf{Средно} & \textbf{Ст. откл.} & \textbf{Декодиране} \\
\midrule
PostgreSQL, текстов & 3638 & 4479 & 3817 & 585 & 1846 \\
PostgreSQL, двоичен & 3115 & 5117 & 3333 & 715 & 250 \\
MySQL, текстов      & 5585 & 5991 & 5647 & 304 & не е отделяно \\
\bottomrule
\end{tabular}
\end{table}

Разликата в декодирането е 7,4 пъти: 1846 срещу 250 микросекунди за 5000 реда по три
декодирани колони, тоест 123 срещу 17 наносекунди на стойност. Причината е в кода: двоичното
декодиране на цяло число е едно зареждане и разместване на байтове, докато текстовото минава
през разбор на десетичен запис, а двоичното декодиране на \code{double} е
\code{std::bit\_cast}, докато текстовото вика \code{strtod}. Общото време също е по-малко за
двоичния формат, но печели \emph{по-малко}, отколкото декодирането спестява: 523 микросекунди
срещу 1596. Клиентът декодира ред, докато сървърът още произвежда следващите, тоест част от
спестеното време освобождава време, което така или иначе щеше да се изчака.

\begin{table}[H]
\centering
\caption{Закъснителност на заемане на връзка от пул, микросекунди}
\label{tab:pool}
\begin{tabular}{L{1.8cm}L{2.2cm}L{2.0cm}L{2.0cm}L{2.0cm}L{2.2cm}}
\toprule
\textbf{Размер} & \textbf{Едновремен.} & \textbf{Медиана} & \textbf{95-и пр.} & \textbf{Ст. откл.} & \textbf{Чакали} \\
\midrule
1 & 1  & 2,1  & 3,5   & 2,9 & 0 от 40 \\
1 & 4  & 1957 & 2213  & 245 & 159 от 160 \\
1 & 16 & 9570 & 10363 & 956 & 639 от 640 \\
2 & 1  & 0,6  & 1,2   & 0,2 & 0 от 40 \\
2 & 4  & 749  & 904   & 121 & 158 от 160 \\
2 & 16 & 5210 & 6384  & 650 & 638 от 640 \\
4 & 1  & 0,5  & 0,7   & 0,2 & 0 от 40 \\
4 & 4  & 0,5  & 1,2   & 0,5 & 0 от 160 \\
4 & 16 & 2343 & 2630  & 297 & 636 от 640 \\
8 & 1  & 0,4  & 0,8   & 0,1 & 0 от 40 \\
8 & 4  & 0,4  & 0,9   & 0,5 & 0 от 160 \\
8 & 16 & 808  & 1028  & 165 & 630 от 640 \\
\bottomrule
\end{tabular}
\end{table}

Когато едновременността не надхвърля размера на пула, заемането струва под три микросекунди и
нито едно заемане не изчаква. Когато я надхвърли, заемането се превръща в чакане на опашка и
цената му е времето за изпълнение на чуждо заявление, тоест хиляди микросекунди вместо части
от една. Разликата между двата режима е три порядъка и е най-голямата, която измерването
откри. При едновременност 16 медианата пада от 9570 микросекунди при пул от една връзка до
808 при пул от осем, тоест осемкратното увеличение намалява чакането 11,8 пъти, повече от
пропорционално: при по-голям пул по-малка част от заявките изобщо стигат до опашката, което
личи и от последната колона. Стандартното отклонение при насищане е около една пета от
медианата или по-малко, тоест величината е чакане от самия пул, а не шум от средата.

Разликите се отнасят към три източника. \textbf{Брой копия}: двоичният формат печели 7,4 пъти
в декодирането, разликата е изцяло в клиента и е точно това, което прякото говорене по
протокол прави наблюдаемо, а именно един ред в Табл.~\ref{tab:formats} вместо скрито
поведение в чужда библиотека. \textbf{Брой съобщения}: кешът спестява 42 на сто, точно по
броене, но съответната разлика във времето е под разсейването и затова не се твърди.
\textbf{Режим на работа}: заемането от пул има два режима с три порядъка разлика, а границата
е едновременност, равна на размера на пула; това е единственият показател, при който изборът
на параметър променя цената качествено.

Четири неща не са измерени. Сравнение с \code{libpq} и \code{libmysqlclient} липсва изцяло,
защото те не са налични на машината; програмата го съобщава и не отчита число, а твърдение за
отношението към доставчиковите библиотеки не се прави. Върховата потребена памет не се
отчита, затова колона за памет липсва във всички таблици. Измерване при мрежа с ненулево
закъснение не е правено, а то е единственият начин да се провери дали спестените 42 на сто
съобщения се превръщат в спестено време. Измерена е закъснителността на заемане от пул, но не
и пропускливостта, тоест броят завършени заявления за единица време.

\TODO{сравнение с \code{libpq} и \code{libmysqlclient} на машина, на която те са инсталирани;
клонът в измервателната програма съществува и чака само средата}

\TODO{измерване на върховата потребена памет}


% ============================================================================
% ЗАКЛЮЧЕНИЕ, без номер
% ============================================================================
% ЗАКЛЮЧЕНИЕ, без номер.
\section*{ЗАКЛЮЧЕНИЕ}
\label{sec:conclusion}
\addcontentsline{toc}{section}{ЗАКЛЮЧЕНИЕ}

Проектът реализира клиентска библиотека на C++20, която говори пряко протоколите на PostgreSQL
и MySQL, с асинхронен пул, кеш на подготвените заявления и декодиране без копие. Изборът е
обоснован в Раздел~\ref{sec:analysis}: прякото говорене прави цената на достъпа наблюдаема.

\textbf{Предимства.} Липсата на междинни слоеве премахва копия, които иначе са скрити, и
измерването го потвърждава: декодирането от двоичен формат е 7,4 пъти по-евтино от текстовото,
а изборът на формат е достъпен само защото клиентът сам изпраща \code{Bind}. Кешът спестява 42
на сто от протоколните съобщения за едно действие. Асинхронният модел не изисква нишка на
връзка: всичко се изпълнява в един цикъл на събитията, без ключове за заключване, а проектът
се сглобява само с компилатор и CMake.

\textbf{Недостатъци и ограничения.} Проектът поема две протоколни реализации и риска от
разминаване при нови версии. Обхватът на протоколите не е пълен, съзнателно. Извън него
остават: за PostgreSQL режимът \code{COPY},
\code{LISTEN}/\code{NOTIFY}, отмяната на изпълняващо се заявление по втора връзка, TLS и
GSSAPI или SSPI; за MySQL двоичният протокол за подготвени заявления, компресията,
\code{caching\_sha2\_password} и отново TLS.
Декодирането обхваща едномерни масиви, а не вложени, и прехвърля отговорност за живота на
данните върху приложението. Без TLS библиотеката е годна в доверена мрежа и не е годна през
публична. Сравнение с \code{libpq} и \code{libmysqlclient} не е направено, защото те не са
налични на машината, и твърдение за отношението към тях не се прави.

\textbf{Инженерна трактовка.} Три пъти измерването поправи очакване. Разширеният цикъл
\emph{не} струва допълнително обръщане при първо изпълнение, защото петте съобщения се
изпращат в един запис. Измерването откри и дефект, който тестовете не бяха уловили: при пул с
размер едно единият работник гладуваше, защото освободената връзка се връщаше в общия списък;
поправката е пряко предаване, покрита с тест. Третото е методическо: случаят с пропуск в кеша
първоначално измерваше и изместване, тоест даваше 15 съобщения и 2 обръщания, вярно за
изпълняваното и грешно спрямо етикета си.

\textbf{По-нататъшна работа.} Сравнение с доставчиковите библиотеки на подходяща машина; TLS;
трети протокол, който би проверил доколко общият слой е общ; съвместимост с CLI~\cite{sql_cli}.


% ============================================================================
% Списък на използваните литературни източници
% ============================================================================
\newpage
\phantomsection
\addcontentsline{toc}{section}{ЛИТЕРАТУРА}
\bibliographystyle{ieeetr}
\bibliography{references}

% ============================================================================
% Приложения
% ============================================================================
% ПРИЛОЖЕНИЕ, без номер. Само указател към изходния текст, без вмъкване на код.
\section*{ПРИЛОЖЕНИЕ}
\label{sec:appendix}
\addcontentsline{toc}{section}{ПРИЛОЖЕНИЕ}

Изходният текст не се препечатва тук. Хранилището е публично:
\url{https://github.com/Alex-Tsvetanov/conduit}. Табл.~\ref{tab:sources} дава указател към
него; редовете са преброени с \code{wc -l}.

\begin{table}[H]
\centering
\caption{Указател към изходния текст на проекта}
\label{tab:sources}
\begin{tabular}{L{4.4cm}L{7.6cm}L{1.6cm}}
\toprule
\textbf{Път} & \textbf{Какво съдържа} & \textbf{Редове} \\
\midrule
\code{bytes.hpp}, \code{task.hpp}, \code{net.*}, \code{tls.*} & буфери, \code{task<T>}, сокет, TLS и цикъл на събитията & 1509 \\
\code{crypto.*}, \code{scram.*} & MD5, SHA-1, SHA-256, HMAC, PBKDF2, SCRAM-SHA-256 & 470 \\
\code{codec.hpp}, \code{value.*}, \code{trace.*} & договорът между кодеците, декодерите, броене на съобщенията & 612 \\
\code{lru\_cache.hpp}, \code{row\_map.hpp}, \code{pool.hpp} & кеш с изместване, типово изображение, асинхронен пул & 496 \\
\code{pg/}, \code{src/pg\_*.cpp} & кодек и връзка за PostgreSQL, двата цикъла, TLS & 1176 \\
\code{mysql/}, \code{src/mysql\_*.cpp} & кодек и връзка за MySQL, текст протокол, TLS & 852 \\
\midrule
\code{tests/} & 91 модулни и 14 интеграционни случая, собствена рамка & 2406 \\
\code{benchmarks/bench.cpp} & измервателната програма, флаг \code{--raw} & 613 \\
\code{benchmarks/symmetric\_transfer\_check.cpp} & проверка дали симетричният трансфер расте по стек & 85 \\
\code{examples/demo.cpp} & показна програма с отпечатан протоколен обмен & 286 \\
\code{docker-compose.yml} & PostgreSQL 16.4 и MySQL 8.0.39, фиксирани версии, SSL & 87 \\
\bottomrule
\end{tabular}
\end{table}

Схемата на измервателната таблица се създава от самата измервателна програма и е част от
нейния изходен текст. Суровите проби от всяко повторение се записват при подаден флаг
\code{--raw} във файла \code{raw\_samples.csv}; пускът, от който са таблиците в
Раздел~\ref{sec:results}, произведе 5660 записа с полета \code{case}, \code{repetition} и
\code{microseconds}.


\end{document}
